\section{The agent: how a diagnosis happens}
\label{sec:agent}

The agent is an LLM (GPT-5.4, hosted on Azure) in a loop with seven read-only
investigation tools over one recorded incident. It is told only that ``an anomaly was
detected'' in a known time window --- nothing about what or where --- and it must end
by committing to a verdict: \emph{which component} is the root cause, \emph{what kind
of fault} it is (from a fixed menu of 13 types), the evidence, and a confidence.

\paragraph{The tools.} Each tool compares the healthy baseline window against the
incident window and returns a compact summary --- never raw data:
\begin{itemize}
  \item \textbf{traces}: per-service request latency percentiles;
  \item \textbf{topology}: which service calls which, and which of those hops slowed
  down --- including calls into components that emit no traces of their own
  (databases, queues), recovered from the callers' side;
  \item \textbf{logs}: how each service's \emph{error rate changed}, plus error
  messages that are new since the incident began (long-standing noise is flagged as
  such, so it cannot masquerade as a cause);
  \item \textbf{metrics}: resource signals per container ranked by how much they
  moved, host-wide signals, and ``limit'' evidence (a container throttled by its CPU
  cap, or pinned at a memory ceiling);
  \item \textbf{kernel}: the derived kernel views described above --- changed activity,
  plain-English digests, and the wait explanation;
  \item \textbf{source code}: search and read the application's own code (useful for
  checking what a timeout or error message actually does);
  \item plus a service list for orientation.
\end{itemize}

\paragraph{The context builder.} Before the agent starts, a deterministic component
runs the survey pass itself --- what changed, where, on every data source --- and
hands the agent a one-page evidence briefing. This ``brief'' turned out to be the
single best accuracy-per-dollar improvement we made: same or better accuracy at
roughly half the tokens and 40\% fewer investigation steps.

\paragraph{The skill library.} On top of the generic method, we wrote one short
investigation guide per fault type --- what that problem looks like in evidence, how
to confirm it, and how to tell it apart from look-alikes. A selector reads the
evidence briefing and either picks the best-matching guide or abstains, in which case
the agent proceeds from first principles. Guides are single text files, deliberately
written so that anyone can add one for a problem they care about.

\paragraph{Keeping the evaluation honest.} Injected-fault datasets can quietly hand
the model the answer: recording names that contain the fault type, containers named
after the fault, label files. We strip or pseudonymise everything of that kind before
the model sees it, and an automated checker re-scans every experiment's full record to
certify nothing leaked. This mattered: with the giveaways present, the agent looked
far better than it really was (it was effectively reading the answer off identifiers).
All numbers in this report are from the honest setup, and every diagnosis ships with
its complete, replayable record.
